\section{Question and result}
\label{sec:question}

Resonant-mass gravitational antennas are historically described through elastic eigenmodes, absorption cross sections, effective coupling, directivity, and frequency-integrated response \cite{Hirakawa1976,PaikWagoner1976,Lobo1995,Aguiar2011}. Material-response dispersion relations, gravitational-wave sum rules, and complete generator--receiver calculations provide related cumulative-response viewpoints \cite{Srivastava2003,Rudenko2003}. Generic source--receiver wave-channel bounds and passive $H_2$ machinery are also established \cite{Miller2000,Molesky2020,BarasBrockett1975,Opmeer2013}. Classical matching and antenna theory likewise contain gain--bandwidth restrictions \cite{Fano1950,Chu1948,Harrington1960}. The generic principle that passive resonant structure trades peak response against bandwidth is therefore not the contribution claimed here.

The question is narrower and gravity-specific: when \emph{both} compact gravitational endpoints are allowed arbitrary passive linear-harmonic complexity within a specified retained modal sector, can cumulative coherent transfer be reduced to endpoint mass-distribution resources and the propagating transverse-traceless (TT) channel?

Let $\omega_0$ be the absolute carrier frequency, $\nu$ the complex-envelope detuning, and
\begin{equation}
\omega(\nu)=\omega_0+\nu .
\end{equation}
The measured detuning band is $\mathcal B_\nu$, with $B/\omega_0\ll1$, and
\begin{equation}
\omega_-=\inf_{\nu\in\mathcal B_\nu}\omega(\nu)>0 .
\label{eq:omegamin}
\end{equation}
Let $a_A$ and $a_B$ denote characteristic endpoint radii. The compact separated-wave regime requires the complete measured band to satisfy
\begin{equation}
\frac{\omega_+a_A}{c}\ll1,\qquad
\frac{\omega_+a_B}{c}\ll1,\qquad
\frac{\omega_-R}{c}\gg1 ,
\label{eq:geometryscope}
\end{equation}
where $\omega_+=\sup_{\nu\in\mathcal B_\nu}\omega(\nu)$.

Let $T(\nu)$ map selected energy-normalized local input channels at source $A$ to selected local output channels at receiver $B$, with all unobserved passive ports retained in the physical model. Define
\begin{equation}
\boxed{
\Gamma_{\rm coh}
\equiv
\frac{1}{2\pi}
\int_{\mathcal B_\nu}\dd\nu\,
\Tr[T^\dagger(\nu)T(\nu)] .
}
\label{eq:gamma}
\end{equation}
For a scalar channel the integrand is power transmissivity. Thus $\Gamma_{\rm coh}$ is the area under the transmissivity spectrum; equivalently it is a band-limited squared $H_2$ norm of the selected transfer block. It has units of inverse time. Without a noise model and signaling protocol it is neither an information capacity nor a bit rate.

Choose the separation direction $\hat R$ and define, about each endpoint center of mass,
\begin{align}
I_2&=\int_V\rho r^2\dd^3x,\\
I_{\hat R}&=\int_V\rho\left[r^2-(\hat R\cdot\bm x)^2\right]\dd^3x,\\
Z_{\hat R}&=\int_V\rho(\hat R\cdot\bm x)^2\dd^3x ,
\end{align}
so $I_2=I_{\hat R}+Z_{\hat R}$. Here $I_{\hat R}$ is the conventional moment of inertia about the line joining the endpoints.

\begin{theorem}[Two-ended passive gravitational spectral area]
Consider separated compact nonrelativistic linear-harmonic source and receiver systems in weak mass-quadrupole gravity satisfying Eq.~\eqref{eq:geometryscope}. At each endpoint let the complex-envelope model contain a finite or countably infinite bounded-port Markov modal sector with retained physical modal frequencies
\begin{equation}
\omega_n\le\Omega .
\label{eq:retainedsector}
\end{equation}
Then the far-zone coefficient obeys
\begin{equation}
\boxed{
\limsup_{R\rightarrow\infty}
R^2\Gamma_{\rm coh}
\le
\frac{5G\Omega^4}{4c^3\omega_-^2}
\min(I_{\hat R,A},I_{\hat R,B}) .
}
\label{eq:asymptotic-headline}
\end{equation}
Repeated passive returns between the same two compact endpoints do not change this leading coefficient.
\end{theorem}

For a retained carrier-scale sector, $\Omega=\omega_0[1+O(B/\omega_0)]$ and $\omega_-=\omega_0[1+O(B/\omega_0)]$. Equation~\eqref{eq:asymptotic-headline} therefore gives the leading narrowband form
\begin{equation}
\boxed{
\Gamma_{\rm coh}
\lesssim
\frac{5G\omega_0^2}{4c^3R^2}
\min(I_{\hat R,A},I_{\hat R,B}) .
}
\label{eq:headline}
\end{equation}
The $\lesssim$ symbol denotes only this final narrowband and far-zone reduction. Section~\ref{sec:propagation} gives a finite-band, finite-$kR$ inequality within the outgoing compact-quadrupole angular model before either approximation is taken.

The retained-sector condition is substantive. The on-shell modal linewidth in Sec.~\ref{sec:resource} is evaluated at the mode's own resonance frequency and supports a near-resonant Markov description only for the retained sector. Far-detuned higher modes cannot consistently be imported into that same model merely by assigning them their on-shell $\omega_n^4$ linewidths. Their low-frequency tails require a frequency-dependent elastic/radiative treatment or additional material regularity. Completeness by itself controls an unweighted quadrupole sum, not an unrestricted fourth frequency moment.

The physical content is therefore narrower than a general detector limit but stronger than a peak-response statement: within the declared compact passive modal class, passive quality factor, retained resonance count, unitary internal hybridization, and same-endpoint recurrence can reshape the transfer spectrum without increasing its cumulative area beyond a geometry-resolved endpoint resource and the TT propagation channel.
